\documentclass[10pt,a4paper]{amsart}
\usepackage[margin=19mm]{geometry}
\usepackage{amsmath,amssymb,mathtools}
\usepackage[scr=rsfs,cal=euler]{mathalfa}
\usepackage{xfrac}
\usepackage[unicode,hidelinks,bookmarks=false]{hyperref}
\newcommand{\ie}{i.e.\ }
\newcommand{\cf}{cf.\ }
\newcommand{\resp}{respectively\ }
\newcommand{\Subsection}{\subsection}
\providecommand{\nobreakdash}{\nobreak\hbox{-}}
\setlength{\emergencystretch}{2em}
\multlinegap0pt
\usepackage{amssymb}
\usepackage{enumerate}
\usepackage{xcolor}
\usepackage{float}
\newtheorem{theorem}{Theorem}
\newtheorem{lemma}{Lemma}
\newcommand{\fN}{\mathfrak N}
\title{The reduced ring of the $RO(C_2)$-graded $C_2$-equivariant stable stems}
\author{Eva Belmont, Zhouli Xu, Shangjie Zhang}
\date{Source: arXiv:2211.04611v1 (CC BY 4.0)}
\begin{document}
\maketitle

\noindent\emph{Source and licence.} The reduced ring of the $RO(C_2)$-graded $C_2$-equivariant stable stems. Version arXiv:2211.04611v1 (CC BY 4.0), distributed by arXiv under the Creative Commons Attribution 4.0 International licence, http://creativecommons.org/licenses/by/4.0/, as recorded in the arXiv licence metadata for this exact version and shown on the abstract page https://arxiv.org/abs/2211.04611v1. Full licence notice: \texttt{licenses/arxiv-2211.04611-CC-BY-4.0.txt}.\par\smallskip

\noindent\textit{Editorial scope.} Counted unit: the article's lemma lem:2-to-eta (source0.tex lines 758--762). The article's complete original proof of it is delivered with it (source0.tex lines 763--778, one \texttt{proof} environment of 16 lines). No mathematics is written, repaired or re-derived: the statement and the proof are the article's own lines, in the article's own notation, and no retained line is substituted or re-flowed.  Retained uncounted as the source-local context the proof needs: lines 405--407; lines 749--751.  Nothing was omitted from any retained span, so the package carries no omission record.  No retained line contains a citation, so the wrapper carries no bibliography and no bibliography entry.  No retained line contains a figure environment, an image-inclusion command or drawing code, so no image is delivered and the package declares no asset dependency.  Editorial formatting only, in addition: an AMS \texttt{amsart} preamble that restates the article's own declarations and macros used by the retained lines, namely the environments \texttt{theorem}, \texttt{lemma} on separate counters, together with the standard packages those lines need.  The retained environments print in this document as Theorem 1, Lemma 1 in order of appearance, whereas the article prints the same items with the numbering of its own full document.  The complete native source of the article as distributed in the arXiv e-print for this exact version is bundled unchanged as \texttt{sources/133212/source0.tex}.
\begin{theorem}\label{thm: BS, BGH}
An element $\alpha\in \pi_\star^{C_2}$ is nilpotent if and only if both its geometric fixed points $\Phi^{C_2}(\alpha)\in \pi_*^s$ and its underlying non-equivariant map $\Phi^e(\alpha)\in \pi_*^s$ are nilpotent.
\end{theorem}

\begin{lemma}\label{lem:2eta}
We have $2\eta = \rho \eta^2$ and $2\rho = \rho^2 \eta$ in $\pi_\star^{C_2}$.
\end{lemma}

\begin{lemma}\label{lem:2-to-eta}
For $0\leq j\leq i$,
we have that $\eta^i$ is divisible by $2^j$ in $\pi_\star^{C_2}/\fN$ if and only if
$\eta^{i-j}$ is divisible by $\rho^j$ in $\pi_\star^{C_2}/\fN$.
\end{lemma}

\begin{proof}
First we show that $\pi_\star^{C_2}/\fN$ is uniquely $\eta$-divisible and
$\rho$-divisible in stems $\neq 0$; i.e., for homogeneous elements $x,y\in
\pi_{s,w}^{C_2}/\fN$ with $s\neq 0$,
if $\eta x = \eta y$ then $x = y$, and similarly for $\rho$.
If $\rho x = \rho y$, then $x-y$ is $\rho$-torsion, which implies
$\Phi^{C_2}(x-y)=0$. Since the stem is nonzero, $\Phi^e(x-y)=0$, and so $x-y$ is
nilpotent by Theorem \ref{thm: BS, BGH}.
If $\eta x = \eta y$, then $\rho^2 \eta x = \rho^2 \eta y$, and using Lemma
\ref{lem:2eta} we have $2\rho x = 2\rho y$. Since $\pi_\star^{C_2}/\fN$ is
torsion-free, we have $\rho x = \rho y$, reducing to the previous case.

Suppose $\eta^i = 2^j x$ modulo $\fN$. Then $\eta^i \rho = 2^j \rho x = (\rho
\eta)^j \rho x$ by Lemma \ref{lem:2eta} and $\eta^{i-j} = \rho^j x$ by $\eta$-
and $\rho$-divisibility. This argument is reversible, establishing the converse.
\end{proof}

\end{document}
